\originalpreface
晶体中电子的运动受周期性约束，而电子彼此之间的作用又使单粒子图像出现边界。学习这两个层次，需要把能带、实验响应、低能有效理论和量子多体态接起来。本册读者应已掌握本科量子力学中的算符、微扰与角动量，以及基本热力学、统计分布和傅里叶变换；二次量子化、推迟响应和 Berry 相位在正文重新建立。

第一编以 Patrick A. Lee 的 \href{https://ocw.mit.edu/courses/8-511-theory-of-solids-i-fall-2004/}{8.511，Fall 2004} 为范围，覆盖公开的24讲主题和11套作业。该课原PDF主要是授课摘要，正文的完整计算链由编者补写，不能视为课堂逐字译稿。第二编接续 Lee 的 \href{https://ocw.mit.edu/courses/8-512-theory-of-solids-ii-spring-2009/}{8.512，Spring 2009}，覆盖25讲主题、12套当学期作业及页面所列2004年补充作业；公开讲义1--10还须署名 Mark Spencer Rudner。课程页讲次与个别PDF内部标题存在错位，来源表保留两种定位。第三编依 Xiao-Gang Wen 的 \href{https://ocw.mit.edu/courses/8-513-modern-quantum-many-body-physics-for-condensed-matter-systems-fall-2021/}{8.513，Fall 2021}，四组幻灯片覆盖1--24讲，作业页列11套题。

8.511成绩由作业70\%与一次测验30\%组成，测验原卷未公开；8.512按12次作业评分；8.513要求每周作业与约四页学期论文，未提供论文范本或官方标准解。课堂摘要引用 Marder、Ashcroft--Mermin 等商业书的部分，只保留真实定位，不把书本引用当成已经取得的题面。作业答案若无公开官方解，明确标记为编者独立补写。

原作者为上述教师及获署名的讲义作者，课程平台为 Massachusetts Institute of Technology OpenCourseWare。中文翻译、结构重编、推导补充和可编辑TikZ图均为本项目改编。基于OCW的部分遵守\href{https://creativecommons.org/licenses/by-nc-sa/4.0/}{CC BY-NC-SA 4.0}：署名、非商业、相同方式共享；不表示MIT或原作者认可本项目。8.513幻灯片中明确标注第三方例外的图片不进入本稿，也不随源码包分发。逐件链接、页数、哈希、许可例外及未取得项目见卷末和学科目录中的来源清单。
